% TeX-ready caption and results text for the public-v661 main figure.
\textbf{Structural credit routing and focal shunting in a frozen MICrONS v661 sensitivity cohort.}
\textbf{A}, Composition of the frozen 47-cell cohort by excitatory cell class.
\textbf{B}, Cell-mean modeled field capture across feedback-channel budgets for a dense PCA oracle, morphology-defined paths, random paths, depth bins and ancestry shuffles ($n=47$ cells through eight channels; $n=46$ at 16 channels). Twenty perturbation streams were averaged within each cell and condition; bands are 95\% cell-bootstrap confidence intervals.
\textbf{C}, Cell-level capture at eight channels.
\textbf{D}, Eight-channel wiring-capture trade-off; morphology paths used 8.1\% of dense-feedback wiring.
\textbf{E}, Within-cell morphology advantage over the three structural controls.
\textbf{F}, At dose 1, localization after a focal shunt and first-order-current-matched additive perturbation ($n=45$ cells, 235 focal sites).
\textbf{G}, Localization for the true descendant relation and relation templates reassigned within cell ($n=40$ cells, 230 focal sites).
\textbf{H}, Direct E/I-label coverage and eligible focal-site count by cell class. Diamonds and intervals in \textbf{C,E--G} show cell means and 95\% cell-bootstrap confidence intervals. Cells, not focal sites or streams, are the independent units. The cohort excludes the original eight cells by stable nucleus identifier but comes from the same MICrONS mouse. It uses historical static v661 reconstructions and direct presynaptic coarse E/I calls only (11,024 annotations before spatial mapping, 10,516 mapped contacts; 4.71\% of incoming synapses). These are structural and passive-network model tests, not measurements of in vivo learning.

\paragraph{Disjoint-cell sensitivity analysis in public MICrONS v661 reconstructions.}
We froze a 55-cell V1 excitatory cohort before computing sensitivity outcomes and excluded the eight discovery cells by stable nucleus identifier. All 47 remaining v661 skeletons and postsynaptic meshworks were available from the public static release. Twenty perturbation streams were averaged within each cell. At eight feedback channels, morphology-defined paths captured 75.6\% of weighted modeled field energy using 8.1\% of dense-feedback wiring, compared with 83.9\% for the dense PCA oracle. Morphology capture exceeded random paths by 45.3 percentage points (95\% CI 41.8--48.8), depth bins by 31.5 points (28.1--34.9) and ancestry shuffles by 42.5 points (40.0--45.0); all three contrasts were positive in 47/47 cells. Across 235 eligible sites in 45 cells, focal shunting increased descendant-selective gradient localization relative to the first-order-current-matched additive control by 0.138 (0.116--0.162; 45/45 cells). The true relation also exceeded reassigned relations by 0.118 (0.094--0.144; 39 positive and one numerical tie among 40 cells). Because this analysis uses a historical release from the same animal and direct E/I labels cover 4.71\% of incoming synapses (11,024 before mapping; 10,516 mapped), it supports robustness across disjoint reconstructed cells, not independent-animal generalization or an in vivo learning claim.
